%% ============================================================
%%  TahalilAI — Doctor Directory & Recommendation System
%%  Engineering Development Report
%%  Date: March 2026
%% ============================================================
\documentclass[12pt, a4paper]{article}

% ---------- Packages ----------
\usepackage[T1]{fontenc}
\usepackage[utf8]{inputenc}
\usepackage[english]{babel}
\usepackage{lmodern}
\usepackage{microtype}
\usepackage{geometry}
\usepackage{graphicx}
\usepackage{xcolor}
\usepackage{booktabs}
\usepackage{longtable}
\usepackage{array}
\usepackage{tabularx}
\usepackage{multirow}
\usepackage{enumitem}
\usepackage{fancyhdr}
\usepackage{titlesec}
\usepackage{hyperref}
\usepackage{listings}
\usepackage{tcolorbox}
\usepackage{tikz}
\usepackage{pgfplots}
\usepackage{amsmath}
\usepackage{float}
\usepackage{caption}
\usepackage{subcaption}
\usepackage{mdframed}
\usepackage{fontawesome5}
\usepackage{setspace}
\usepackage{parskip}
\usepackage{colortbl}
\usepackage{makecell}

\pgfplotsset{compat=1.18}
\usetikzlibrary{shapes.geometric, arrows.meta, positioning, shadows,
                fit, backgrounds, decorations.pathreplacing, calc}

\geometry{
  top=2.5cm, bottom=2.5cm,
  left=2.8cm, right=2.8cm,
  headheight=15pt
}

% ---------- Colour Palette ----------
\definecolor{primary}{HTML}{1A3C5E}
\definecolor{accent}{HTML}{2980B9}
\definecolor{success}{HTML}{27AE60}
\definecolor{warning}{HTML}{F39C12}
\definecolor{danger}{HTML}{E74C3C}
\definecolor{purple}{HTML}{8E44AD}
\definecolor{teal}{HTML}{16A085}
\definecolor{lightgray}{HTML}{F5F7FA}
\definecolor{midgray}{HTML}{BDC3C7}
\definecolor{codebg}{HTML}{F0F4F8}
\definecolor{codefg}{HTML}{2C3E50}
\definecolor{dbcolor}{HTML}{2ECC71}
\definecolor{apicolor}{HTML}{3498DB}
\definecolor{frontcolor}{HTML}{9B59B6}
\definecolor{aicolor}{HTML}{E67E22}
\definecolor{etlcolor}{HTML}{1ABC9C}

% ---------- Hyperref ----------
\hypersetup{
  colorlinks = true,
  linkcolor  = accent,
  urlcolor   = accent,
  citecolor  = accent,
  pdftitle   = {TahalilAI Doctor Feature Report},
  pdfauthor  = {TahalilAI Development Team}
}

% ---------- Section Formatting ----------
\titleformat{\section}
  {\Large\bfseries\color{primary}}
  {\thesection}{1em}{}
  [\color{midgray}\titlerule]

\titleformat{\subsection}
  {\large\bfseries\color{accent}}
  {\thesubsection}{1em}{}

\titleformat{\subsubsection}
  {\normalsize\bfseries\color{primary}}
  {\thesubsubsection}{1em}{}

% ---------- Header / Footer ----------
\pagestyle{fancy}
\fancyhf{}
\fancyhead[L]{\small\color{primary}\textbf{TahalilAI} — Doctor Directory \& Recommendation Report}
\fancyhead[R]{\small\color{midgray}March 2026}
\fancyfoot[C]{\small\color{midgray}\thepage}
\renewcommand{\headrulewidth}{0.4pt}
\renewcommand{\headrule}{\hbox to\headwidth{\color{midgray}\leaders\hrule height
  \headrulewidth\hfill}}

% ---------- Code Listing Style ----------
\lstdefinestyle{pythonstyle}{
  backgroundcolor   = \color{codebg},
  basicstyle        = \ttfamily\footnotesize\color{codefg},
  keywordstyle      = \bfseries\color{accent},
  commentstyle      = \itshape\color{success},
  stringstyle       = \color{danger},
  numberstyle       = \tiny\color{midgray},
  numbers           = left,
  numbersep         = 8pt,
  breaklines        = true,
  breakatwhitespace = true,
  showstringspaces  = false,
  frame             = single,
  rulecolor         = \color{midgray},
  tabsize           = 4,
  language          = Python,
  xleftmargin       = 10pt,
  xrightmargin      = 5pt,
  aboveskip         = 8pt,
  belowskip         = 8pt,
  captionpos        = b,
}

\lstdefinestyle{tscode}{
  backgroundcolor   = \color{codebg},
  basicstyle        = \ttfamily\footnotesize\color{codefg},
  keywordstyle      = \bfseries\color{purple},
  commentstyle      = \itshape\color{success},
  stringstyle       = \color{danger},
  numberstyle       = \tiny\color{midgray},
  numbers           = left,
  numbersep         = 8pt,
  breaklines        = true,
  showstringspaces  = false,
  frame             = single,
  rulecolor         = \color{midgray},
  tabsize           = 2,
  xleftmargin       = 10pt,
  aboveskip         = 8pt,
  belowskip         = 8pt,
  captionpos        = b,
}
\lstset{style=pythonstyle}

% ---------- tcolorbox Styles ----------
\tcbuselibrary{breakable, skins, listings}

\newtcolorbox{infobox}[1][]{
  enhanced, breakable,
  colback=lightgray, colframe=accent,
  fonttitle=\bfseries\color{primary},
  left=6pt, right=6pt, top=4pt, bottom=4pt,
  #1
}

\newtcolorbox{successbox}[1][]{
  enhanced, breakable,
  colback=green!5!white, colframe=success,
  fonttitle=\bfseries\color{success},
  left=6pt, right=6pt, top=4pt, bottom=4pt,
  #1
}

\newtcolorbox{warningbox}[1][]{
  enhanced, breakable,
  colback=yellow!5!white, colframe=warning,
  fonttitle=\bfseries\color{warning},
  left=6pt, right=6pt, top=4pt, bottom=4pt,
  #1
}

\newtcolorbox{notebox}[1][]{
  enhanced, breakable,
  colback=purple!5!white, colframe=purple,
  fonttitle=\bfseries\color{purple},
  left=6pt, right=6pt, top=4pt, bottom=4pt,
  #1
}

% ---------- Custom Commands ----------
\newcommand{\file}[1]{\texttt{\color{accent}#1}}
\newcommand{\func}[1]{\texttt{\color{danger}#1}}
\newcommand{\envvar}[1]{\texttt{\color{warning}#1}}
\newcommand{\col}[1]{\texttt{\color{teal}#1}}
\newcommand{\badge}[2]{\colorbox{#1}{\color{white}\small\textbf{\ #2\ }}}
\newcommand{\endpoint}[2]{\texttt{\badge{apicolor}{#1}\ \color{primary}#2}}

% ============================================================
\begin{document}
% ============================================================

% ---------- Title Page ----------
\begin{titlepage}
  \begin{tikzpicture}[remember picture, overlay]
    \fill[primary] (current page.north west)
      rectangle ([yshift=-5.5cm] current page.north east);
    \fill[accent] ([yshift=-5.5cm] current page.north west)
      rectangle ([yshift=-6cm] current page.north east);
  \end{tikzpicture}

  \vspace*{1.5cm}
  \begin{center}
    {\color{white}\fontsize{28}{34}\selectfont\bfseries TahalilAI}\\[0.4cm]
    {\color{white!80!gray}\fontsize{14}{18}\selectfont Engineering Development Report}\\[0.2cm]
    {\color{accent!60!white}\fontsize{11}{14}\selectfont
      Doctor Directory, Recommendation System \& AI-Powered Specialist Matching}
  \end{center}

  \vspace{4.5cm}

  \begin{center}
    \begin{tcolorbox}[
      enhanced, width=0.88\textwidth,
      colback=lightgray, colframe=primary,
      arc=4pt, boxrule=1pt,
      left=14pt, right=14pt, top=10pt, bottom=10pt
    ]
      \begin{tabular}{@{}ll}
        \textbf{Project}    & TahalilAI — AI-Powered Medical Analysis Platform \\[4pt]
        \textbf{Module}     & Doctor Directory \& Recommendation System \\[4pt]
        \textbf{Authors}    & TahalilAI Development Team \\[4pt]
        \textbf{Date}       & March 4, 2026 \\[4pt]
        \textbf{Version}    & 1.0 \\[4pt]
        \textbf{Records}    & 9,220 doctors across Morocco \\[4pt]
        \textbf{Status}     & \badge{success}{Released} \\
      \end{tabular}
    \end{tcolorbox}
  \end{center}

  \vfill
  \begin{center}
    \color{midgray}\small
    Faculty of Sciences and Techniques — Computer Science\\
    School Project — Academic Year 2025--2026
  \end{center}
\end{titlepage}

% ---------- Table of Contents ----------
\newpage
\tableofcontents
\newpage

% ============================================================
\section{Executive Summary}
% ============================================================

This report documents the design, implementation, and testing of the
\textbf{Doctor Directory and AI-Powered Recommendation System} added to the
TahalilAI medical analysis platform.
The feature enables patients to discover qualified doctors in Morocco,
receive specialist recommendations based on their AI-generated analysis
results, and interact with a rich doctor directory through the web interface.

The implementation spans eight distinct phases:

\begin{enumerate}[leftmargin=*, itemsep=4pt]
  \item \textbf{CSV Data Cleanup} — A multi-step ETL pipeline that normalises
        raw doctor data from 131 unique speciality variants to 78 canonical
        names, cleans phone numbers, merges duplicate records, and
        removes 1,407 placeholder phone entries.
  \item \textbf{SQLite Database} — SQLAlchemy ORM model with 17 fields,
        3 individual indexes, and a composite index for performant queries.
  \item \textbf{Backend API} — Six RESTful endpoints (list, filters, detail,
        recommendation, geolocation) registered on the FastAPI application.
  \item \textbf{AI Recommendation Service} — Gemini-powered speciality
        extraction with structured JSON output and local fallback.
  \item \textbf{Frontend Directory Page} — Full-featured React page with
        debounced search, multi-filter dropdowns, smart pagination,
        and a detail modal.
  \item \textbf{Frontend Recommendation Components} — Three reusable
        components integrating seamlessly with the results page and
        displaying urgency-aware doctor cards.
  \item \textbf{Optional Enrichment Scripts} — Batch geocoding and
        Google Places rating enrichment pipelines.
  \item \textbf{Comprehensive Test Suite} — 35 unit tests covering
        ORM CRUD, query logic, and AI recommendation behaviour.
\end{enumerate}

\begin{infobox}[title={Module Metrics at a Glance}]
  \begin{tabular}{lrrr}
    \toprule
    \textbf{Layer} & \textbf{Files} & \textbf{Lines} & \textbf{Tests} \\
    \midrule
    ETL Scripts (data pipeline) & 4 & $\approx$350 & — \\
    Database \& Models          & 2 & $\approx$170 & 17 \\
    Backend API Router          & 1 & $\approx$160 & — \\
    Recommender Service         & 1 & $\approx$180 & 18 \\
    Pydantic Schemas            & 1 & $\approx$60  & — \\
    Frontend Components         & 4 & $\approx$650 & — \\
    \midrule
    \textbf{Total}              & \textbf{13} & $\approx$\textbf{1,570} & \textbf{35} \\
    \bottomrule
  \end{tabular}
\end{infobox}

% ============================================================
\section{Full-Stack Architecture}
% ============================================================

\begin{figure}[H]
\centering
\begin{tikzpicture}[
  node distance=0.5cm and 2cm,
  every node/.style={font=\small},
  layer/.style={rectangle, rounded corners=3pt, minimum width=2.6cm,
                minimum height=0.8cm, draw, thick, align=center},
  arrow/.style={-{Stealth[length=6pt]}, thick},
  dashed arrow/.style={-{Stealth[length=6pt]}, dashed, color=midgray}
]

%% --- Row 0: Data Sources ---
\node[layer, fill=etlcolor!20, draw=etlcolor] (csv) {Raw CSV\\(9,220 rows)};
\node[layer, fill=etlcolor!20, draw=etlcolor, right=1.6cm of csv] (clean) {Cleaned CSV\\(deduplicated)};
\node[layer, fill=warning!20, draw=warning, right=1.6cm of clean] (gmaps) {Google Maps\\API};

%% --- Row 1: Database ---
\node[layer, fill=dbcolor!20, draw=dbcolor, below=1.4cm of clean] (db) {SQLite DB\\(doctors table)};

%% --- Row 2: Backend ---
\node[layer, fill=apicolor!20, draw=apicolor, below left=1.4cm and -0.2cm of db] (router) {FastAPI Router\\(/doctors)};
\node[layer, fill=aicolor!20, draw=aicolor, below right=1.4cm and -0.2cm of db] (recommender) {Recommender\\Service};

%% --- Row 2.5: AI ---
\node[layer, fill=aicolor!15, draw=aicolor, right=2cm of recommender] (gemini) {Gemini\\Flash};

%% --- Row 3: Frontend ---
\node[layer, fill=frontcolor!20, draw=frontcolor, below=2.8cm of router] (page) {Doctors Page\\(directory)};
\node[layer, fill=frontcolor!20, draw=frontcolor, right=2cm of page] (recdoctors) {Recommended\\Doctors};

%% Arrows
\draw[arrow, color=etlcolor] (csv) -- (clean);
\draw[arrow, color=etlcolor] (clean) -- (db);
\draw[dashed arrow]           (gmaps) -- node[above, font=\tiny]{geocode/rate} (db);
\draw[arrow, color=apicolor]  (db) -- (router);
\draw[arrow, color=apicolor]  (db) -- (recommender);
\draw[arrow, color=apicolor]  (router) -- (page);
\draw[arrow, color=frontcolor](recommender) -- (recdoctors);
\draw[dashed arrow]           (recommender) -- node[above, font=\tiny]{prompt} (gemini);
\draw[dashed arrow]           (gemini) -| (recommender);

%% Layer labels
\node[font=\tiny\itshape\color{midgray}, above=0.05cm of csv] {Phase 1: ETL};
\node[font=\tiny\itshape\color{midgray}, above=0.05cm of db] {Phase 2: Database};
\node[font=\tiny\itshape\color{midgray}, below=0.05cm of page] {Phase 4–5: Frontend};

\end{tikzpicture}
\caption{End-to-end data flow through the Doctor Directory feature}
\label{fig:architecture}
\end{figure}

% ============================================================
\section{Phase 1 — CSV Data Cleaning Pipeline}
% ============================================================

\subsection{Overview}

The raw source file \file{backend/doctors\_data.csv} contains 9,220 rows scraped
from Moroccan medical directory websites. The script
\file{backend/scripts/clean\_doctors\_csv.py} transforms this raw data into a
production-ready \file{doctors\_data\_clean.csv} through twelve processing stages.

\subsection{Processing Stages}

\begin{longtable}{p{0.4cm} p{3.8cm} p{9.2cm}}
  \toprule
  \textbf{\#} & \textbf{Stage} & \textbf{Description} \\
  \midrule
  1  & Load (BOM-safe)        & UTF-8-BOM encoding stripping via \texttt{encoding="utf-8-sig"} \\[3pt]
  2  & Drop unused columns    & Removes \texttt{email}, \texttt{neighborhood}, \texttt{coordinates} \\[3pt]
  3  & Extract title          & Strips "Dr." / "Dr " prefix from name; stores in \col{title} column \\[3pt]
  4  & Normalise phones       & Converts all Moroccan numbers to \texttt{+212XXXXXXXXX} E.164 format \\[3pt]
  5  & Clear placeholders     & Removes junk values like \texttt{+212522000000} (1,407 entries cleared) \\[3pt]
  6  & Extract primary spec.  & Sets \col{primary\_speciality} from first item in speciality list \\[3pt]
  7  & Normalise cities       & Maps city variants: Témara $\to$ Temara, Laâyoune $\to$ Laayoune, etc. \\[3pt]
  8  & Normalise specialities & Maps 131 raw variants to 78 canonical French names \\[3pt]
  9  & Clean descriptions     & Removes \texttt{||} / \texttt{|} separators, collapses whitespace \\[3pt]
  10 & Clean image URLs       & Replaces 21 known Cloudinary placeholder URLs with empty strings \\[3pt]
  11 & Drop sparse columns    & Drops \col{timetable} ($\approx$90.6\% empty) and \col{cabinet\_name} ($\approx$85.2\% empty) \\[3pt]
  12 & Deduplicate            & Groups by (name, city), keeps "richest" row by field completeness score \\
  \bottomrule
  \caption{CSV data cleaning pipeline stages}
  \label{tab:etl_stages}
\end{longtable}

\subsection{Key Normalisation Examples}

\begin{infobox}[title={Speciality Normalisation Samples}]
  \begin{tabular}{ll}
    \textbf{Raw Value} & \textbf{Canonical Name} \\
    \hline
    Médecine générale & Médecin généraliste \\
    Chirurgien viscéral & Chirurgien digestif \\
    Chirurgien maxillo-facial & Chirurgien général \\
    Médecin interniste & Médecin interniste \\
  \end{tabular}
\end{infobox}

\subsection{Deduplication Strategy}

The \func{\_row\_richness(row)} helper scores each row by: (1) the number of
non-empty fields it contains, and (2) the length of the \col{description} field.
When multiple rows share the same (name, city) pair, only the highest-scoring
row is retained. This strategy preserves the most informative record without
manual intervention.

\subsection{Output Schema}

The cleaned CSV contains the following ordered columns:

\begin{center}
  \small\texttt{title, name, primary\_speciality, specialities, phone,
  address, city, description, languages, image\_url, profile\_url, source\_site}
\end{center}

% ============================================================
\section{Phase 2 — Database Layer}
% ============================================================

\subsection{Database Configuration (\file{database.py})}

The database layer is built with \textbf{SQLAlchemy 2.0+} and a
\textbf{SQLite} backend located at \file{backend/tahalilai.db}.

\begin{itemize}[leftmargin=*, itemsep=3pt]
  \item \textbf{Engine}: Created with
        \texttt{connect\_args=\{"check\_same\_thread": False\}} for
        FastAPI thread-safety compatibility.
  \item \textbf{Session factory}: \texttt{SessionLocal} with
        \texttt{autocommit=False, autoflush=False}.
  \item \textbf{Dependency injection}: The \func{get\_db()} generator
        yields a scoped session per HTTP request and closes it in a
        \texttt{finally} block, preventing connection leaks.
\end{itemize}

\subsection{Doctor ORM Model (\file{models.py})}

\begin{longtable}{p{3.8cm} p{2cm} p{1.8cm} p{5.4cm}}
  \toprule
  \textbf{Column} & \textbf{Type} & \textbf{Nullable} & \textbf{Notes} \\
  \midrule
  \col{id}                  & Integer & No  & Primary key, auto-increment \\
  \col{title}               & String  & Yes & "Dr", empty string default \\
  \col{name}                & String  & No  & Individual index \texttt{ix\_name} \\
  \col{primary\_speciality} & String  & No  & Individual index \texttt{ix\_speciality} \\
  \col{specialities}        & String  & Yes & Comma-separated list \\
  \col{phone}               & String  & Yes & E.164 format (\texttt{+212...}) \\
  \col{address}             & String  & Yes & Physical office address \\
  \col{city}                & String  & No  & Individual index \texttt{ix\_city} \\
  \col{description}         & String  & Yes & Free-text biography \\
  \col{languages}           & String  & Yes & Spoken languages \\
  \col{image\_url}          & String  & Yes & Profile photo URL \\
  \col{profile\_url}        & String  & Yes & External directory link \\
  \col{source\_site}        & String  & Yes & Origin data source \\
  \col{latitude}            & Float   & Yes & WGS84 (geocoded, optional) \\
  \col{longitude}           & Float   & Yes & WGS84 (geocoded, optional) \\
  \col{google\_rating}      & Float   & Yes & Google Places rating (0.0–5.0) \\
  \col{google\_review\_count}& Integer& Yes & Number of Google reviews \\
  \bottomrule
  \caption{Doctor model column definitions}
  \label{tab:doctor_columns}
\end{longtable}

\begin{infobox}[title={Composite Index}]
  A composite index \texttt{ix\_doctors\_city\_speciality} is defined on
  \texttt{(city, primary\_speciality)}. This accelerates the most common
  query pattern: \emph{``find doctors of speciality X in city Y''},
  which is the core of the recommendation endpoint.
\end{infobox}

\subsection{Database Seeder (\file{seed\_db.py})}

The \func{seed()} function in \file{backend/scripts/seed\_db.py}:

\begin{enumerate}[leftmargin=*, itemsep=3pt]
  \item Calls \texttt{Base.metadata.create\_all()} to create all tables.
  \item Checks if the \texttt{doctors} table already contains rows
        (idempotent — skips if already seeded).
  \item Reads \file{doctors\_data\_clean.csv} and instantiates
        \texttt{Doctor} objects for each row.
  \item Inserts all records in a single \texttt{bulk\_save\_objects()} call,
        loading \textbf{9,220 doctors} in one batch commit.
\end{enumerate}

% ============================================================
\section{Phase 3 — Backend API}
% ============================================================

\subsection{Router Configuration}

The doctor endpoints are organised in
\file{backend/src/tahalilai/routers/doctors.py} and registered with:

\begin{itemize}[leftmargin=*]
  \item \textbf{Prefix:} \texttt{/doctors}
  \item \textbf{Tags:} \texttt{["doctors"]}
  \item \textbf{Dependency:} \texttt{db: Session = Depends(get\_db)}
        applied per route.
\end{itemize}

\subsection{Endpoint Reference}

\subsubsection{\endpoint{GET}{/doctors} — List \& Search}

\begin{longtable}{p{3cm} p{2.5cm} p{7cm}}
  \toprule
  \textbf{Parameter} & \textbf{Type} & \textbf{Description} \\
  \midrule
  \texttt{city}       & \texttt{str | None} & Case-insensitive city filter \\
  \texttt{speciality} & \texttt{str | None} & Case-insensitive speciality filter \\
  \texttt{q}          & \texttt{str | None} & Full-text search: name, address, description \\
  \texttt{page}       & \texttt{int = 1}    & Page number ($\geq 1$) \\
  \texttt{page\_size} & \texttt{int = 20}   & Results per page ($1 \leq n \leq 100$) \\
  \bottomrule
\end{longtable}

Returns \texttt{DoctorListResponse} with \texttt{doctors[]}, \texttt{total},
\texttt{page}, and \texttt{page\_size}.
All string filters use SQLAlchemy \texttt{ilike()} for case-insensitive
partial matching.

\subsubsection{\endpoint{GET}{/doctors/cities} — City Aggregation}

Returns \texttt{list[CityCount]} — cities ordered by doctor count descending.
Query: \texttt{SELECT city, COUNT(id) GROUP BY city ORDER BY COUNT(id) DESC}.

\subsubsection{\endpoint{GET}{/doctors/specialities} — Speciality Aggregation}

Returns \texttt{list[SpecialityCount]} — specialities ordered by doctor count
descending.

\subsubsection{\endpoint{GET}{/doctors/\{doctor\_id\}} — Single Doctor}

Returns a single \texttt{DoctorResponse}. Raises \texttt{HTTP 404} with a
descriptive message if the ID is not found.

\subsubsection{\endpoint{GET}{/doctors/recommend} — AI Recommendations}

\begin{longtable}{p{3cm} p{2.5cm} p{7cm}}
  \toprule
  \textbf{Parameter} & \textbf{Type} & \textbf{Description} \\
  \midrule
  \texttt{specialities} & \texttt{str} (required) & Comma-separated speciality list \\
  \texttt{city}         & \texttt{str | None}   & Optional city preference \\
  \texttt{limit}        & \texttt{int = 5}       & Max doctors per speciality ($1 \leq n \leq 20$) \\
  \bottomrule
\end{longtable}

\textbf{Algorithm:} For each received speciality, the endpoint queries doctors
filtered by \texttt{primary\_speciality}. When a city is provided, doctors
from that city are prepended to fill the result set first; remaining slots
are filled from any city. All results are deduplicated by doctor ID before
returning.

\subsubsection{\endpoint{GET}{/doctors/nearby} — Geolocation Search}

\begin{longtable}{p{3cm} p{2.5cm} p{7cm}}
  \toprule
  \textbf{Parameter} & \textbf{Type} & \textbf{Description} \\
  \midrule
  \texttt{lat}          & \texttt{float} (required) & User latitude (WGS84) \\
  \texttt{lng}          & \texttt{float} (required) & User longitude (WGS84) \\
  \texttt{speciality}   & \texttt{str | None}       & Optional speciality filter \\
  \texttt{radius\_km}   & \texttt{float = 10.0}     & Search radius ($0.1 \leq r \leq 100$) km \\
  \texttt{limit}        & \texttt{int = 10}         & Max results ($1 \leq n \leq 50$) \\
  \bottomrule
\end{longtable}

\textbf{Algorithm:} Fetches all doctors with non-null coordinates, applies
the \textbf{Haversine formula} (Earth radius $R = 6371$ km) to compute
great-circle distances, filters by radius, sorts by distance ascending,
and appends a \texttt{distance\_km} field (rounded to 1 decimal) to each
result.

\begin{lstlisting}[caption={Haversine distance calculation}, language=Python]
import math

def _haversine(lat1: float, lng1: float, lat2: float, lng2: float) -> float:
    R = 6371.0
    phi1, phi2 = math.radians(lat1), math.radians(lat2)
    dphi = math.radians(lat2 - lat1)
    dlam = math.radians(lng2 - lng1)
    a = math.sin(dphi / 2)**2 + math.cos(phi1) * math.cos(phi2) * math.sin(dlam / 2)**2
    return R * 2 * math.atan2(math.sqrt(a), math.sqrt(1 - a))
\end{lstlisting}

\subsection{Pydantic Schemas (\file{schemas.py})}

\begin{longtable}{p{4cm} p{9.5cm}}
  \toprule
  \textbf{Schema} & \textbf{Fields} \\
  \midrule
  \texttt{DoctorResponse}      & 17 fields mirroring the ORM model; \texttt{from\_attributes=True} for ORM serialisation \\[4pt]
  \texttt{DoctorListResponse}  & \texttt{doctors: list[DoctorResponse]}, \texttt{total}, \texttt{page}, \texttt{page\_size} \\[4pt]
  \texttt{CityCount}           & \texttt{city: str}, \texttt{count: int} \\[4pt]
  \texttt{SpecialityCount}     & \texttt{speciality: str}, \texttt{count: int} \\
  \bottomrule
  \caption{Doctor-related Pydantic schemas}
\end{longtable}

% ============================================================
\section{Phase 4 — AI Recommendation Service}
% ============================================================

\subsection{Overview}

The recommender service (\file{services/recommender.py}) bridges the
AI analysis pipeline and the doctor database. After an analysis job
completes, the service uses \textbf{Gemini Flash} to read the analysis
text and extract the most relevant medical specialities the patient
should consult.

\subsection{Canonical Speciality List}

The module maintains a list of \textbf{36 canonical French medical
speciality names} (\texttt{CANONICAL\_SPECIALITIES}) covering all major
disciplines represented in the Moroccan doctor database. This list is
injected into the Gemini prompt so that extracted speciality names are
guaranteed to be valid database values.

\begin{infobox}[title={Sample Canonical Specialities (subset)}]
  Médecin généraliste, Cardiologue, Dermatologue, Endocrinologue,\\
  Gastro-entérologue, Gynécologue, Neurologue, Ophtalmologue,\\
  Pédiatre, Pneumologue, Psychiatre, Radiologue, Rhumatologue, Urologue,\\
  Diabétologue, Nutritionniste, Kinésithérapeute, Oncologue, \ldots
\end{infobox}

\subsection{Gemini Integration}

\subsubsection{\func{extract\_recommended\_specialities(analysis)}}

\begin{itemize}[leftmargin=*, itemsep=3pt]
  \item \textbf{Input:} Full analysis text string (truncated to 3,000
        characters in the prompt to stay within token budget).
  \item \textbf{Prompt:} Instructs Gemini to return a strict JSON object
        with two keys: \texttt{"specialities"} (list from
        \texttt{CANONICAL\_SPECIALITIES}) and \texttt{"urgency"}
        (\texttt{"routine"|"soon"|"urgent"}).
  \item \textbf{Model:} Gemini Flash, \texttt{temperature=0.1},
        \texttt{max\_output\_tokens=256}.
  \item \textbf{Post-processing:} Strips markdown code fences
        (\texttt{```json...```}) before JSON parsing.
  \item \textbf{Default:} Falls back to
        \texttt{\{"specialities": ["Médecin généraliste"], "urgency": "routine"\}}
        on any failure — empty API key, network error, or malformed JSON.
\end{itemize}

\begin{figure}[H]
  \centering
  \begin{tikzpicture}[
    node distance=0.9cm,
    box/.style={rectangle, rounded corners=3pt, minimum width=5cm,
                minimum height=0.75cm, draw, thick, align=center, font=\small},
    decision/.style={diamond, draw, thick, aspect=2, font=\small,
                     fill=lightgray, minimum width=4.5cm},
    arrow/.style={-{Stealth}, thick, color=primary}
  ]
    \node[box, fill=primary!15, draw=primary] (input)  {Analysis Text};
    \node[box, fill=lightgray, draw=accent, below of=input] (check)  {API key present \& Gemini available?};
    \node[box, fill=aicolor!15, draw=aicolor, below=1.2cm of check] (gemini) {Gemini Flash API};
    \node[box, fill=lightgray, draw=accent, below of=gemini] (parse)  {Strip fences → parse JSON};
    \node[box, fill=success!15, draw=success, below of=parse] (ok)    {Return \{specialities, urgency\}};
    \node[box, fill=danger!10, draw=danger, right=3.5cm of check] (default) {Default: généraliste\\+ routine};

    \draw[arrow] (input)  -- (check);
    \draw[arrow] (check)  -- node[left, font=\tiny]{yes} (gemini);
    \draw[arrow] (gemini) -- (parse);
    \draw[arrow] (parse)  -- (ok);
    \draw[arrow, color=danger] (check) -- node[above, font=\tiny]{no / error} (default);
    \draw[arrow, color=danger, dashed] (gemini.east) -- ++(0.8,0) |- node[right, font=\tiny]{exception} (default);
  \end{tikzpicture}
  \caption{Speciality extraction decision flow}
  \label{fig:extraction_flow}
\end{figure}

\subsubsection{\func{find\_recommended\_doctors(specialities, city, db, limit\_per\_speciality)}}

\begin{itemize}[leftmargin=*, itemsep=3pt]
  \item For each speciality in the input list, queries the database
        using \texttt{ilike(\%speciality\%)}.
  \item When a city is provided, city-matching doctors are prepended;
        remaining slots are filled from any city.
  \item Deduplicates across all specialities using a \texttt{seen\_ids} set.
  \item Returns a list of dicts with keys:
        \texttt{id, title, name, speciality, phone, address, city,
        image\_url, profile\_url}.
\end{itemize}

% ============================================================
\section{Phase 5 \& 6 — Frontend Components}
% ============================================================

\subsection{Doctors Directory Page (\file{app/doctors/page.tsx})}

The \textbf{Doctors Page} is a Next.js client component providing a
full-featured interface for browsing the doctor database.

\begin{longtable}{p{3.5cm} p{9.5cm}}
  \toprule
  \textbf{Feature} & \textbf{Implementation} \\
  \midrule
  Search           & Debounced (300 ms) text input querying name, address, and description \\[3pt]
  City filter      & Dropdown populated from \texttt{GET /doctors/cities} (sorted by count) \\[3pt]
  Speciality filter& Dropdown populated from \texttt{GET /doctors/specialities} \\[3pt]
  Clear filters    & Single button resets city, speciality, search, and page simultaneously \\[3pt]
  Grid layout      & 1 / 2 / 3 columns responsive (mobile / tablet / desktop) \\[3pt]
  Pagination       & Smart pagination showing up to 5 page buttons; Previous/Next at boundaries \\[3pt]
  Doctor cards     & Rendered via \texttt{DoctorCard} components; click opens detail modal \\[3pt]
  Loading state    & Spinner replaces grid during API fetch \\[3pt]
  Empty state      & Friendly "No results" message with suggestion \\
  \bottomrule
  \caption{Doctors page features}
\end{longtable}

\textbf{State management:} 10 \texttt{useState} hooks manage independent
concerns (doctors, filters, pagination, modal). Three \texttt{useEffect}
hooks handle: (1) initial filter option loading, (2) search debouncing,
and (3) re-fetching doctors when any filter or page changes.

\begin{lstlisting}[language=Python, caption={Search debounce effect (TypeScript)},
  style=tscode]
// 300 ms debounce prevents an API call on every keystroke
useEffect(() => {
  const handler = setTimeout(() => setDebouncedQuery(searchQuery), 300);
  return () => clearTimeout(handler);
}, [searchQuery]);
\end{lstlisting}

\subsection{DoctorCard Component (\file{components/DoctorCard.tsx})}

A compact, clickable card displaying the most relevant doctor information
at a glance.

\begin{itemize}[leftmargin=*, itemsep=3pt]
  \item \textbf{Avatar:} Profile image if available; else a coloured
        monogram badge (doctor initials).
  \item \textbf{Name \& Speciality:} Truncated name in bold; primary
        speciality in accent colour.
  \item \textbf{City \& Distance:} Location icon with city name;
        optional \texttt{distance\_km} label (used by the nearby endpoint).
  \item \textbf{Rating:} Star icon with Google rating and review count
        when available.
  \item \textbf{Phone:} Displayed as small grey text when present.
  \item \textbf{Interaction:} \texttt{cursor-pointer} with hover shadow
        and ring; triggers \texttt{onClick} prop callback.
\end{itemize}

The \texttt{Doctor} TypeScript interface exported by this component is
reused by \texttt{DoctorDetailModal} and \texttt{RecommendedDoctors}.

\subsection{DoctorDetailModal Component (\file{components/DoctorDetailModal.tsx})}

A full-screen modal overlay providing complete doctor information and
direct-action buttons.

\begin{itemize}[leftmargin=*, itemsep=3pt]
  \item \textbf{Header:} Gradient banner (indigo $\to$ blue) with
        avatar, name, and primary speciality.
  \item \textbf{Body (scrollable):} City, address, phone, languages,
        Google rating, and description (truncated to 500 characters).
  \item \textbf{Footer actions:}
    \begin{itemize}
      \item \textbf{Call} button — \texttt{<a href="tel:...">} for
            one-tap dialling.
      \item \textbf{View Profile} button — opens external directory
            link in a new tab (rendered only when
            \texttt{profile\_url} is present).
    \end{itemize}
  \item \textbf{Dismissal:} Backdrop click or close button (top-right)
        fires the \texttt{onClose} prop callback.
\end{itemize}

\subsection{RecommendedDoctors Component (\file{components/RecommendedDoctors.tsx})}

This component is embedded in the analysis results page and surfaces
AI-matched doctors with urgency-aware styling.

\begin{longtable}{p{2.5cm} p{3.5cm} p{7cm}}
  \toprule
  \textbf{Urgency} & \textbf{Gradient} & \textbf{Banner Message} \\
  \midrule
  \texttt{urgent}  & Red (red-600 $\to$ rose-600) & "Urgent — Please consult as soon as possible" \\
  \texttt{soon}    & Amber (amber-600 $\to$ orange-600) & "Recommended — Schedule a visit soon" \\
  routine          & Green (emerald-600 $\to$ teal-600) & "Routine — For your regular follow-up" \\
  \bottomrule
  \caption{Urgency-based visual variants in RecommendedDoctors}
\end{longtable}

\textbf{Data mapping:} The component converts the
\texttt{RecommendedDoctor} interface (returned by the recommendation API,
with a \texttt{speciality} key) to the \texttt{Doctor} interface expected
by \texttt{DoctorCard} (which uses \texttt{primary\_speciality}), filling
all other optional fields with empty strings.

% ============================================================
\section{Phase 7 — Optional Enrichment Scripts}
% ============================================================

\subsection{Geocoding Script (\file{scripts/geocode\_doctors.py})}

\begin{itemize}[leftmargin=*, itemsep=3pt]
  \item Queries doctors where \col{latitude} \texttt{IS NULL} and
        \texttt{address} is non-empty.
  \item For each doctor, calls the \textbf{Google Maps Geocoding API}
        with the query \texttt{"\{address\}, \{city\}, Morocco"}.
  \item Updates \col{latitude} and \col{longitude} columns in the database.
  \item Batches commits every 50 records; includes 0.1 s sleep to
        respect rate limits.
  \item Requires \envvar{GOOGLE\_MAPS\_API\_KEY}.
\end{itemize}

\begin{warningbox}[title={Cost Estimate}]
  Google Geocoding API charges approximately \$5 per 1,000 requests.
  Geocoding all 9,220 doctors once costs roughly \$46.
  The free tier (up to 40,000 requests/month) covers a full run at no cost.
\end{warningbox}

\subsection{Rating Enrichment Script (\file{scripts/enrich\_ratings.py})}

\begin{itemize}[leftmargin=*, itemsep=3pt]
  \item Queries doctors where \col{google\_rating} \texttt{IS NULL}.
  \item For each doctor, calls the \textbf{Google Places Text Search API}
        with the query \texttt{"Dr \{name\} \{city\} Morocco"}.
  \item Updates \col{google\_rating} and
        \col{google\_review\_count} columns.
  \item Supports an optional city argument (\texttt{sys.argv[1]}) for
        selective enrichment.
  \item Batches commits every 50 records; includes 0.2 s rate-limit sleep.
\end{itemize}

\begin{warningbox}[title={Cost Estimate}]
  Google Places Text Search charges approximately \$32 per 1,000 requests.
  A full 9,220-doctor enrichment run costs roughly \$295.
  Selective enrichment (top cities only) is recommended for cost control.
\end{warningbox}

% ============================================================
\section{Test Suite}
% ============================================================

\subsection{Testing Philosophy}

All tests use \textbf{in-memory SQLite} databases (created fresh per test
function or class) to ensure complete isolation from the production database.
Gemini API calls are patched with \texttt{@patch} decorators so tests are
fully deterministic and run without any credentials or network access.

\subsection{Test Module: Database Models}

\textbf{File:} \file{backend/tests/unit/test\_models.py} \quad
(243 lines, 17 test cases)

\begin{longtable}{p{3.3cm} p{3.5cm} p{6.5cm}}
  \toprule
  \textbf{Class} & \textbf{Test} & \textbf{Assertion} \\
  \midrule
  \multirow{4}{*}{\texttt{TestDoctorSchema}}
    & \texttt{test\_table\_name}       & \texttt{\_\_tablename\_\_ == "doctors"} \\
    & \texttt{test\_table\_created}    & Table exists after \texttt{create\_all()} \\
    & \texttt{test\_columns\_exist}    & All 17 expected columns present \\
    & \texttt{test\_indexes\_exist}    & Composite index \texttt{ix\_doctors\_city\_speciality} confirmed \\
  \midrule
  \multirow{8}{*}{\texttt{TestDoctorCRUD}}
    & \texttt{test\_create}            & Insert returns \texttt{id > 0} \\
    & \texttt{test\_read}              & Fetch by id, validate fields \\
    & \texttt{test\_update}            & Modify fields, re-fetch, assert changes \\
    & \texttt{test\_delete}            & Delete, confirm query returns \texttt{None} \\
    & \texttt{test\_default\_values}   & Empty strings for title/phone/address, \texttt{None} for coords/ratings \\
    & \texttt{test\_nullable\_fields}  & lat/lng/rating accept \texttt{None} \\
    & \texttt{test\_geocoding\_fields} & Lat/lng stored and retrieved with $\delta < 0.001$ \\
    & \texttt{test\_rating\_fields}    & Rating and review count stored correctly \\
  \midrule
  \multirow{5}{*}{\texttt{TestDoctorQueries}}
    & \texttt{test\_filter\_city}       & \texttt{ilike} city returns 2 records \\
    & \texttt{test\_filter\_speciality} & \texttt{ilike} speciality returns 2 records \\
    & \texttt{test\_filter\_combined}   & Combined filters return 1 record \\
    & \texttt{test\_search\_by\_name}   & Wildcard name search works \\
    & \texttt{test\_count\_by\_city}    & \texttt{GROUP BY} city produces correct counts \\
  \bottomrule
  \caption{Database model test cases (17 total)}
  \label{tab:model_tests}
\end{longtable}

\subsection{Test Module: Recommender Service}

\textbf{File:} \file{backend/tests/unit/test\_recommender.py} \quad
(254 lines, 18 test cases)

\begin{longtable}{p{3.8cm} p{3cm} p{6.5cm}}
  \toprule
  \textbf{Class} & \textbf{Test} & \textbf{Assertion} \\
  \midrule
  \multirow{4}{*}{\makecell[l]{\texttt{TestCanonical}\\\texttt{Specialities}}}
    & \texttt{test\_not\_empty}         & List has $\geq 1$ entries \\
    & \texttt{test\_contains\_gen.}     & "Médecin généraliste" in list \\
    & \texttt{test\_contains\_card.}    & "Cardiologue" in list \\
    & \texttt{test\_no\_duplicates}     & All entries unique \\
  \midrule
  \multirow{7}{*}{\makecell[l]{\texttt{TestExtract}\\\texttt{Recommended}}}
    & \texttt{test\_default\_no\_gemini}  & Returns généraliste/routine when Gemini absent \\
    & \texttt{test\_default\_no\_key}     & Returns default when API key empty \\
    & \texttt{test\_parses\_response}     & Parses valid JSON specialities + urgency \\
    & \texttt{test\_markdown\_fence}      & Strips \texttt{```json...```} correctly \\
    & \texttt{test\_api\_exception}       & API error → default \\
    & \texttt{test\_invalid\_json}        & Malformed JSON → default \\
    & \texttt{test\_truncates\_long}      & Analysis $> 3000$ chars truncated in prompt \\
  \midrule
  \multirow{8}{*}{\makecell[l]{\texttt{TestFind}\\\texttt{Recommended}}}
    & \texttt{test\_finds\_by\_spec.}     & Returns matching doctors \\
    & \texttt{test\_prioritizes\_city}    & City results appear first \\
    & \texttt{test\_fills\_other\_cities} & Fills slots from other cities \\
    & \texttt{test\_multiple\_specs.}     & Returns results for each speciality \\
    & \texttt{test\_no\_duplicates}       & Each doctor appears at most once \\
    & \texttt{test\_empty\_list}          & Empty specialities → empty results \\
    & \texttt{test\_unknown\_spec.}       & Non-existent speciality → \texttt{[]} \\
    & \texttt{test\_respects\_limit}      & $\leq \texttt{limit\_per\_speciality}$ returned \\
  \bottomrule
  \caption{Recommender service test cases (18 total)}
  \label{tab:recommender_tests}
\end{longtable}

\subsection{Test Coverage Summary}

\begin{figure}[H]
  \centering
  \begin{tikzpicture}
    \begin{axis}[
      ybar, bar width=1.6cm,
      width=10cm, height=7cm,
      xlabel={Test Module}, ylabel={Test Cases},
      symbolic x coords={test\_models, test\_recommender},
      xtick=data,
      ymin=0, ymax=24,
      ymajorgrids=true,
      grid style=dashed,
      nodes near coords,
      nodes near coords align={above},
      every node near coord/.style={font=\small\bfseries},
      axis line style={color=primary},
      tick style={color=midgray},
      label style={color=primary},
    ]
      \addplot[fill=dbcolor!70, draw=dbcolor]
        coordinates {(test\_models, 17)};
      \addplot[fill=aicolor!70, draw=aicolor]
        coordinates {(test\_recommender, 18)};
    \end{axis}
  \end{tikzpicture}
  \caption{Unit test distribution across doctor-feature modules (total: 35 tests)}
  \label{fig:test_bar}
\end{figure}

% ============================================================
\section{Configuration Reference}
% ============================================================

\begin{longtable}{p{5cm} p{3cm} p{5.5cm}}
  \toprule
  \textbf{Variable} & \textbf{Required} & \textbf{Purpose} \\
  \midrule
  \envvar{GEMINI\_API\_KEY}      & \badge{warning}{Recommended} & Gemini Flash for speciality extraction \\[3pt]
  \envvar{GOOGLE\_MAPS\_API\_KEY}& \badge{success}{Optional} & Geocoding and Places rating enrichment \\[3pt]
  \envvar{DATABASE\_URL}         & No (auto-derived) & Override default SQLite path \\
  \bottomrule
  \caption{Doctor feature environment variables}
  \label{tab:config}
\end{longtable}

\subsection{Running the Test Suite}

\begin{lstlisting}[language=bash, style=pythonstyle,
  caption={Executing doctor feature unit tests}]
# From the backend directory
cd backend

# Run all doctor-related unit tests
pytest tests/unit/test_models.py tests/unit/test_recommender.py -v

# Run with coverage report
pytest tests/unit/test_models.py tests/unit/test_recommender.py \
  --cov=tahalilai.models \
  --cov=tahalilai.services.recommender \
  --cov-report=term-missing

# Seed the database (run once after cleaning CSV)
python scripts/clean_doctors_csv.py
python scripts/seed_db.py

# Optional: geocode and enrich ratings
python scripts/geocode_doctors.py       # requires GOOGLE_MAPS_API_KEY
python scripts/enrich_ratings.py        # requires GOOGLE_MAPS_API_KEY
\end{lstlisting}

% ============================================================
\section{Design Decisions}
% ============================================================

\subsection{SQLite over PostgreSQL}

SQLite was chosen for its zero-configuration setup and the fully
self-contained nature of the school project. The SQLAlchemy ORM layer
ensures that migrating to PostgreSQL requires only a one-line change
in the connection URL, with no modifications to model or query code.

\subsection{Haversine in Python vs. PostGIS}

The geolocation endpoint computes distances in Python rather than
delegating to a spatial database extension. For the current scale
(9,220 records, all pre-filtered to non-null coordinates), this
in-memory approach is sufficient and avoids the PostGIS dependency.

\subsection{French Canonical Specialities}

All speciality names are stored and queried in \textbf{French},
the primary language of the Moroccan medical community.
The recommender injects the full canonical list into the Gemini
prompt, ensuring extracted names are always valid database tokens.

\subsection{Graceful Recommendation Fallback}

The Gemini dependency is fully optional. When the API key is absent
or the service is unavailable, the recommendation pipeline silently
falls back to "Médecin généraliste" with urgency "routine" — ensuring
the results page always displays a useful suggestion rather than
failing silently.

% ============================================================
\section{Summary and Conclusion}
% ============================================================

This report has documented the complete end-to-end implementation of the
TahalilAI Doctor Directory and Recommendation System, comprising:

\begin{enumerate}[leftmargin=*, itemsep=5pt]
  \item A \textbf{12-stage ETL pipeline} that transforms raw scraped data
        into a clean, deduplicated, and fully normalised doctor database
        of 9,220 records.
  \item A \textbf{SQLAlchemy ORM} layer with a 17-column Doctor model,
        3 individual indexes, and a composite city/speciality index
        optimised for the recommendation query pattern.
  \item A \textbf{6-endpoint FastAPI router} providing list, search,
        filter, recommendation, geolocation, and detail operations
        with full Pydantic validation.
  \item A \textbf{Gemini-powered recommendation service} that reads
        AI analysis results, extracts relevant French medical specialities,
        and returns a ranked list of matching doctors — with a robust
        local fallback.
  \item A \textbf{complete React frontend} featuring a searchable
        directory page, a reusable doctor card, a full-detail modal,
        and an urgency-aware recommendation panel integrated into the
        analysis results flow.
  \item A \textbf{35-test unit suite} validating all ORM CRUD operations,
        complex database queries, Gemini response parsing, and recommendation
        logic — entirely offline and credential-free.
\end{enumerate}

\begin{successbox}[title={Deliverables Completed}]
  \begin{itemize}[leftmargin=*, itemsep=3pt, topsep=0pt]
    \item[\checkmark] \file{scripts/clean\_doctors\_csv.py} — 12-stage ETL pipeline
    \item[\checkmark] \file{scripts/seed\_db.py} — Bulk CSV-to-SQLite seeder (9,220 records)
    \item[\checkmark] \file{scripts/geocode\_doctors.py} — Batch Google Maps geocoding
    \item[\checkmark] \file{scripts/enrich\_ratings.py} — Batch Google Places rating enrichment
    \item[\checkmark] \file{database.py} — SQLAlchemy engine + session factory
    \item[\checkmark] \file{models.py} — Doctor ORM model, 17 columns, composite index
    \item[\checkmark] \file{schemas.py} — DoctorResponse, DoctorListResponse, CityCount, SpecialityCount
    \item[\checkmark] \file{routers/doctors.py} — 6 FastAPI endpoints
    \item[\checkmark] \file{services/recommender.py} — Gemini extraction + DB lookup
    \item[\checkmark] \file{app/doctors/page.tsx} — Full directory page
    \item[\checkmark] \file{components/DoctorCard.tsx} — Reusable doctor card
    \item[\checkmark] \file{components/DoctorDetailModal.tsx} — Full detail modal
    \item[\checkmark] \file{components/RecommendedDoctors.tsx} — Urgency-aware recommendations
    \item[\checkmark] \file{tests/unit/test\_models.py} — 17 unit tests
    \item[\checkmark] \file{tests/unit/test\_recommender.py} — 18 unit tests
  \end{itemize}
\end{successbox}

\vspace{1cm}
\begin{center}
  \color{midgray}\small
  \textit{TahalilAI Engineering Development Report — Doctor Directory \& Recommendation System}\\
  \textit{Version 1.0 — March 2026 — Generated for academic review}
\end{center}

\end{document}
